\section{Uložení programu a~dat}\label{sec:hw-architektury}

\subsection{Von Neumannova architektura}

Základní představu o~uspořádání počítače poskytuje \emph{von Neumannova architektura}. Jejím podstatným rysem je, že program i~data jsou uloženy ve~společné paměti. Procesor z~ní načítá strojové instrukce a~stejným způsobem přistupuje také k~hodnotám, které mají být zpracovány. Instrukce určují, jaká operace se má provést, odkud se mají získat operandy a~kam se má uložit výsledek.

Tato myšlenka se někdy označuje jako \emph{princip uloženého programu}. Program není pevně zabudován do zapojení počítače, ale je uložen v~paměti podobně jako ostatní data. Změnou obsahu paměti tak můžeme změnit činnost počítače, aniž bychom museli přestavět jeho hardware. Právě tato univerzálnost je jedním z~důvodů, proč se daný princip stal základem obecných počítačů.

\begin{remark}
    Popis architektury uloženého programu se proslavil v~návrhu počítače EDVAC z~roku 1945,~na němž se podílel John von Neumann. Na vzniku myšlenky se podíleli také J.~Presper Eckert,~John Mauchly a~další členové týmu; tradiční název proto neznamená,~že šlo o~práci jediného autora.
\end{remark}

\begin{figure}[ht]
    \centering
    \begin{tikzpicture}[
    x=1cm,
    y=1cm,
    >=Stealth,
    every node/.style={font=\small},
    block/.style={
        draw,
        thick,
        rounded corners,
        align=center,
        minimum height=1.0cm
    },
    arrow/.style={->, thick},
    both/.style={<->, thick}
]
    \node[
        block,
        fill=blue!10,
        minimum width=3.4cm
    ] (cpu) at (0,1.1) {\textbf{Procesor}};

    \node[
        block,
        fill=green!12,
        minimum width=5.0cm,
        minimum height=1.4cm
    ] (memory) at (0,-1.4) {\textbf{Společná paměť}\\ instrukce i~data};

    \node[
        block,
        fill=orange!15,
        minimum width=2.3cm
    ] (input) at (-5.0,1.1) {Vstup};

    \node[
        block,
        fill=orange!15,
        minimum width=2.3cm
    ] (output) at (5.0,1.1) {Výstup};

    \draw[arrow] (input.east) -- (cpu.west);
    \draw[arrow] (cpu.east) -- (output.west);
    \draw[both] (cpu.south) -- node[right, align=left] {data\\ instrukce} (memory.north);
\end{tikzpicture}

    \caption{Základní uspořádání von Neumannovy architektury.}
    \label{fig:hw-von-neumann}
\end{figure}

Paměť v~tomto modelu sama nerozlišuje, zda uložené bity představují data, nebo instrukce. Rozhoduje až způsob, jakým je procesor použije. Procesor opakovaně načte instrukci z~paměti, zjistí její význam a~provede ji. Pokud instrukce pracuje s~daty uloženými v~paměti, musí procesor obvykle provést další paměťové přenosy. Zjednodušeně tak dostáváme stále se opakující cyklus
\[
    \text{načtení instrukce}
    \longrightarrow
    \text{dekódování}
    \longrightarrow
    \text{provedení}.
\]
Jeho podrobnější podobu rozebereme v~podsekci~\ref{subsec:hw-instrukcni-cyklus}.

Společná cesta pro instrukce i~data přináší jednoduché a~univerzální uspořádání, ale současně může omezovat rychlost. Procesor musí s~pamětí komunikovat při načítání instrukcí i~při čtení a~zápisu dat. Pokud je procesor schopen pracovat rychleji, než mu paměť dodává potřebné informace, část času pouze čeká. Toto omezení se označuje jako \emph{von Neumannovo úzké hrdlo}.

\warningbox{Rychlost počítače neurčuje pouze počet operací, které dokáže procesor provést za sekundu. Stejně důležité je, zda procesor dostává instrukce a~data dostatečně rychle. Právě proto bude později tak významná cache a~celá paměťová hierarchie.}

\subsection{Harvardská architektura}

\emph{Harvardská architektura} používá oddělenou paměť pro instrukce a~oddělenou paměť pro data. Obě paměti mohou mít vlastní adresový prostor i~vlastní přenosové cesty. Procesor pak může současně načítat další instrukci a~přenášet data potřebná pro právě prováděný výpočet. Paměť instrukcí navíc může mít jiné vlastnosti než paměť dat, protože se od ní může očekávat jiný způsob používání.

\begin{figure}[ht]
    \centering
    \begin{tikzpicture}[
    x=1cm,
    y=1cm,
    >=Stealth,
    every node/.style={font=\small},
    title/.style={font=\bfseries},
    block/.style={
        draw,
        thick,
        rounded corners,
        align=center,
        minimum height=0.9cm
    },
    both/.style={<->, thick}
]
    \node[title] at (-3.8,2.1) {Von Neumannova architektura};
    \node[
        block,
        fill=blue!10,
        minimum width=2.8cm
    ] (cpua) at (-3.8,0.9) {Procesor};
    \node[
        block,
        fill=green!12,
        minimum width=4.2cm,
        minimum height=1.25cm
    ] (mema) at (-3.8,-1.2) {Společná paměť\\ instrukce i~data};
    \draw[both] (cpua.south) -- node[right] {\scriptsize instrukce i~data} (mema.north);

    \draw[gray!45, thick] (0,-2.1) -- (0,2.4);

    \node[title] at (3.8,2.1) {Harvardská architektura};
    \node[
        block,
        fill=blue!10,
        minimum width=2.8cm
    ] (cpub) at (3.8,0.9) {Procesor};
    \node[
        block,
        fill=green!12,
        minimum width=2.7cm,
        minimum height=1.25cm
    ] (imem) at (2.0,-1.2) {Paměť\\ instrukcí};
    \node[
        block,
        fill=green!12,
        minimum width=2.7cm,
        minimum height=1.25cm
    ] (dmem) at (5.6,-1.2) {Paměť\\ dat};

    \draw[both] (cpub.south west) -- node[left, align=center] {\scriptsize instrukce} (imem.north);
    \draw[both] (cpub.south east) -- node[right, align=center] {\scriptsize data} (dmem.north);
\end{tikzpicture}

    \caption{Porovnání von Neumannovy a~harvardské architektury.}
    \label{fig:hw-architektury}
\end{figure}

Rozdíl obou přístupů zachycuje obrázek~\ref{fig:hw-architektury}. U~von Neumannovy architektury vede mezi procesorem a~společnou pamětí jedna logická cesta pro instrukce i~data. V~harvardském uspořádání má procesor k~dispozici dvě oddělené cesty. To však neznamená, že každý reálný počítač musí přesně odpovídat jednomu z~těchto dvou schémat.

V~praxi se často používá \emph{modifikovaná harvardská architektura}. Z~pohledu programu může existovat jeden společný adresový prostor, zatímco uvnitř procesoru jsou oddělené například instrukční a~datová cache. Počítač se tedy navenek chová podobně jako von Neumannův, ale v~částech, kde je výhodné souběžné načítání, využívá prvky harvardského uspořádání.

\notebox{Von Neumannova a~harvardská architektura jsou především užitečné modely. Skutečné procesory obsahují mnoho dalších mechanismů a~jejich vnitřní uspořádání bývá kombinací více myšlenek.}
